\subsection{Opis parametrów}
\addcontentsline{toc}{subsection}{Opis parametrów}

\textbf{numberOfMachineTypes} - Ilość typów maszyn (procesów) dostępnych w fabryce

\textbf{numberOfMonths} - Ilość miesięcy uwzględnionych w symulacji

\textbf{numberOfProductsTypes} - Ilość typów produktów

\textbf{numberOfScenarios} - Ilość scenariuszy wygenerowanych do symulacji

\textbf{machines[numberOfMachineTypes]} - Wektor typów maszyn (procesów)

\textbf{months[numberOfMonths]} - Wektor miesięcy symulacji

\textbf{products[numberOfProductsTypes]} - Wektor typów produktów

\textbf{machineCount[numberOfMachineTypes]} - Wektor ilości maszyn danego typu

\textbf{timeToProduce[numberOfMachineTypes][numberOfProductsTypes]} - Macierz czasów produkcji danego produktu na danej maszynie

\textbf{maxProductsInMonth[numberOfMonths][numberOfProductsTypes]} - Macierz maksymalnej ilości produktów, jakie można sprzedać w danym miesiącu

\textbf{numberOfHoursInFactory} - Ilość godzin pracy fabryki w miesiącu

\textbf{mu[numberOfProductsTypes]} - Wektor wartości oczekiwanych rozkładu t-Studenta do generacji scenariuszy

\textbf{sigma[numberOfProductsTypes][numberOfProductsTypes]} - Macierz kowariancji dla rozkładu t-Studenta

\textbf{sellProfit[numberOfScenarios][numberOfProductsTypes]} - Macierz wygenerowanych scenariuszy dochodów ze sprzedaży produktów

\textbf{storageCost} - Koszt trzymania jednej sztuki produktu w magazynie przez miesiąc

\textbf{storageMax[numberOfProductsTypes]} - Wektor maksymalnej pojemności magazynu dla każdego typu produktu

\textbf{storageStart[numberOfProductsTypes]} - Wektor ilości początkowej produktów w magazynie

